\chapter*{Glosario}
\addcontentsline{toc}{chapter}{Glosario}
\markboth{Glosario}{Glosario}

\begin{description}

    \item[API]
    Siglas de \textit{Application Programming Interface}. Interfaz que permite la comunicación entre aplicaciones o componentes software.

    \item[Benchmark]
    Programa o carga de trabajo utilizada para evaluar el comportamiento de un sistema bajo unas condiciones controladas.

    \item[BLAS]
    Siglas de \textit{Basic Linear Algebra Subprograms}. Especificación de rutinas para realizar operaciones básicas de álgebra lineal.

    \item[Campaña experimental]
    Conjunto organizado de ejecuciones definido mediante programas de referencia, configuraciones, perfiles y repeticiones.

    \item[Carga CPU-bound]
    Carga de trabajo cuyo rendimiento está condicionado principalmente por la capacidad de cálculo del procesador.

    \item[Carga memory-bound]
    Carga de trabajo cuyo rendimiento está condicionado principalmente por el acceso a memoria.

    \item[Clúster de computadores]
    Conjunto de computadores interconectados que cooperan para ejecutar trabajos y compartir recursos de cómputo.

    \item[CPU]
    Siglas de \textit{Central Processing Unit}. Componente encargado de ejecutar las instrucciones de los programas.

    \item[CSV]
    Siglas de \textit{Comma-Separated Values}. Formato textual utilizado para almacenar información tabular.

    \item[DGEMM]
    Siglas de \textit{Double-Precision General Matrix Multiplication}. Operación de multiplicación de matrices utilizando números en doble precisión.

    \item[EDP]
    Siglas de \textit{Energy--Delay Product}. Métrica que combina la energía utilizada con el tiempo necesario para completar una ejecución.

    \item[Eficiencia energética]
    Relación entre el trabajo realizado o el rendimiento obtenido y la energía empleada para conseguirlo.

    \item[Energía]
    Magnitud que representa la energía utilizada durante un intervalo temporal. En esta memoria se expresa principalmente en julios.

    \item[Hilo lógico]
    Unidad de ejecución visible para el sistema operativo. Un núcleo físico puede proporcionar más de un hilo lógico mediante SMT.

    \item[HPC]
    Siglas de \textit{High Performance Computing}. Utilización de sistemas de cómputo de altas prestaciones para resolver problemas que requieren una elevada capacidad de procesamiento.

    \item[JSON]
    Siglas de \textit{JavaScript Object Notation}. Formato textual estructurado utilizado para almacenar configuraciones, resultados y metadatos.

    \item[Metadatos]
    Información que describe las condiciones de una ejecución, como el programa utilizado, el número de procesos, el nodo, los parámetros y las marcas temporales.

    \item[NUMA]
    Siglas de \textit{Non-Uniform Memory Access}. Arquitectura en la que el tiempo de acceso a memoria puede depender de la relación entre el procesador y el banco de memoria utilizado.

    \item[Potencia]
    Energía utilizada por unidad de tiempo. En esta memoria se expresa en vatios.

    \item[Prometheus]
    Sistema de monitorización utilizado para recopilar, almacenar y consultar series temporales.

    \item[Pushgateway]
    Componente que permite que trabajos de duración finita publiquen métricas para su posterior recopilación por Prometheus.

    \item[RAPL]
    Siglas de \textit{Running Average Power Limit}. Interfaz de Intel que proporciona contadores energéticos asociados a determinados dominios del procesador.

    \item[Repetibilidad]
    Grado de concordancia entre resultados obtenidos al repetir una ejecución bajo las mismas condiciones experimentales.

    \item[SLURM]
    Gestor de trabajos y recursos utilizado para enviar, planificar y ejecutar cargas en el clúster HPMoon.

    \item[SMT]
    Siglas de \textit{Simultaneous Multithreading}. Técnica que permite que un núcleo físico mantenga más de un hilo lógico de ejecución.

    \item[Throughput]
    Cantidad de trabajo completado por unidad de tiempo.

    \item[Trazabilidad]
    Capacidad para relacionar una ejecución con su configuración, sus muestras originales, sus resultados agregados y sus metadatos.

    \item[Vampire]
    Sistema externo empleado para medir la potencia eléctrica y la energía correspondientes al nodo de cómputo.

    \item[VPS]
    Siglas de \textit{Virtual Private Server}. Servidor virtual utilizado para alojar los servicios de observabilidad del proyecto.

\end{description}

\clearpage